\subsection{DH model}
\label{sec:dh}

The DH frames are assigned from the CuRo6R URDF, with frame $\frameid{0}$ at
\texttt{base\_link}, joint $i$ rotating about $z_{i-1}$, and $\vect{q}=\vect{0}$ matching the
URDF home configuration. We use the standard DH transform
${}^{i-1}\!T_i=R_z(\theta_i)\,T_z(d_i)\,T_x(a_i)\,R_x(\alpha_i)$.
Table~\ref{tab:dh} gives the resulting parameters and the geometric choices
that determine them. Each twist $\alpha_i$ rotates $z_{i-1}$ onto $z_i$ about $+x_i$
by the right-hand rule.
The only fixed joint-angle offset is $\theta_3-q_3=\pi/2$.
The full sequence of frame drawings is given in Appendix~\ref{app:dh-constructions}.

\begin{table}[!ht]
\centering
\caption{One valid DH assignment for CuRo6R. Lengths are in metres; parameter angles are in radians, with degrees used in the explanations.}
\label{tab:dh}
\small
\setlength{\tabcolsep}{5pt}
\renewcommand{\arraystretch}{1.30}
\begin{tabularx}{\linewidth}{@{}c c c c c X@{}}
\toprule
$i$ & $d_i$ & $\theta_i$ & $\alpha_i$ & $a_i$ & Geometric construction at home \\
\midrule
1 & $1$ & $q_1$ & $-\pi/3$ & $0.8$
  & Rise $0.5+0.5$, then move $0.8$ along $x_1=x_0$ to the shoulder.
    Its axis $z_1$ lies $30^\circ$ above $+y_0$, requiring a $-60^\circ$
    turn from $z_0$ about $+x_1$. \\
2 & $1.25$ & $q_2$ & $+\pi/2$ & $2$
  & Move $1.25$ along $z_1$, then $2$ along $x_2$ to reach joint 3.
    Here $x_2=x_1$ and $z_2=-y_1$, so the twist is a positive quarter-turn. \\
3 & $0.75$ & $q_3+\pi/2$ & $+\pi/2$ & $0$
  & Axes 3 and 4 meet $0.75$ along $z_2$, so $a_3=0$.
    Choosing $x_3=y_2$ gives the $+90^\circ$ joint-angle offset;
    $y_2\times z_2=x_2=z_3$ gives the positive quarter-turn twist. \\
4 & $2.5$ & $q_4$ & $-\pi/2$ & $0$
  & The wrist is $1.5+1$ along $z_3$ from $O_3$; intersecting axes give $a_4=0$.
    With $x_4=x_3$, reaching $z_4=y_3$ requires a negative quarter-turn. \\
5 & $0$ & $q_5$ & $+\pi/2$ & $0$
  & Axes 5 and 6 meet at $O_5=O_4$, so $d_5=a_5=0$.
    With $x_5=x_4$, reaching $z_5=-y_4$ requires a positive quarter-turn. \\
6 & $1.5$ & $q_6$ & $0$ & $0$
  & The endpoint $O_6$ lies $1.5$ along $z_5$ from the wrist.
    The last two DH $z$ axes coincide and point the same way, giving $a_6=\alpha_6=0$. \\
\bottomrule
\end{tabularx}
\end{table}

\textbf{Endpoint alignment.}
The endpoint $E$ is \texttt{link\_6}. Its origin coincides with $O_6$, while
$x_E=z_6$, $y_E=x_6$ and $z_E=y_6$. The fixed alignment is therefore appended
to the DH chain to recover the URDF endpoint pose:
\begin{equation}
 {}^0T_E=\left(\prod_{i=1}^{6}{}^{i-1}\!T_i\right){}^6T_E,
 \qquad
 {}^6T_E=\begin{pmatrix}
 0&1&0&0\\0&0&1&0\\1&0&0&0\\0&0&0&1
 \end{pmatrix}.
 \label{eq:dh-tool}
\end{equation}
